\subsection{Complex-Valued SSMs}
\label{sec:method:complex}

Modern SSMs are designed with efficiency as the central goal, motivated by the need to scale to larger models and longer sequences.
For instance, successive architectures have progressively simplified the state-transition matrix: S4~\citep{S4} used complex-valued Normal Plus Low Rank (NPLR) matrices, Mamba~\citep{gu2024mambalineartimesequencemodeling} reduced this to a diagonal of reals, and Mamba-2~\citep{dao2024transformersssmsgeneralizedmodels} further simplified it to a single scaled identity matrix.
Although these simplifications largely maintain language modeling performance, recent works~\citep{merrill2025illusionstatestatespacemodels,sarrof2024expressivecapacitystatespace,grazzi2025unlockingstatetrackinglinearrnns} have shown that the restriction to real, non-negative eigenvalue transitions degrades the capabilities of the model on simple state-tracking tasks---here referring primarily to the solvable-group regime ($\text{TC}^0$) such as parity---which can be solved by a one-layer LSTM.
This limitation, formalized in Theorem 1 of \citep{grazzi2024mambacapableincontextlearning}, arises from restricting the eigenvalues of the transition matrix to real numbers, which cannot represent ``rotational'' hidden state dynamics.
For instance, consider the parity function on binary inputs $\{0,1\}$, defined as $\sum_t x_t \bmod 2$.
This task can be performed using update: $\vh_t = \bR(\pi x_t) \vh_{t-1}$, where $\bR(\cdot)$ is a 2-D rotation matrix. 
Such rotational dynamics cannot be expressed with real eigenvalues. 

\subsubsection{Complex SSM with Exponential-Euler Discretization}
To recover this capability, we begin with complex SSMs~\plaineqref{eq:complex-ssm}, which \emph{are} capable of representing state-tracking dynamics. 
We show that, under discretization (\Cref{prop:var-const}), complex SSMs can be formulated as real SSMs with a \emph{block-diagonal transition matrix composed of $2 \times 2$ rotation matrices} (\Cref{prop:complex-to-real-ssm}). 
We then show that this is equivalent to applying \emph{data-dependent rotary embeddings} on both the input and output projections $\bB, \bC$ respectively.
This result establishes a theoretical connection between complex SSMs and data-dependent RoPE embeddings (\Cref{prop:rope-trick}).
Finally, the ``RoPE trick'' used in \citet{su2023roformerenhancedtransformerrotary} allows for an efficient implementation of complex-valued state-transition matrices with minimal computational overhead compared to real-valued SSMs.

\begin{restatable}[Complex-to-Real SSM Equivalence]{proposition}{PropComplexRealSSM}\label{prop:complex-to-real-ssm}
Consider a complex-valued SSM
\begin{align}
\label{eq:complex-ssm}
    \dot{\vh}(t) &= \mathrm{Diag}\big(A(t) + i \boldsymbol{\theta}(t)\big)\,\vh(t) 
    + \big(\bB(t) + i\hat{\bB}(t)\big)\,x(t), \\ \nonumber
    y(t) &= \mathrm{Re}\!\left(
        \big(\bC(t) + i\hat{\bC}(t)\big)^{\top}\vh(t)
    \right),
\end{align}
where $\vh(t) \in \mathbb{C}^{N/2}$, $\boldsymbol{\theta}(t),\bB(t),\hat{\bB}(t),\bC(t),\hat{\bC}(t)\in\mathbb{R}^{N/2}$, and $x(t),A(t)\in\mathbb{R}$.
Under exponential-Euler discretization, this system is equivalent to a real-valued SSM
\begin{align}
\label{eq:real-ssm}
    \vh_t &= e^{\Delta_t A_t}\,\bR_t\,\vh_{t-1} + \Delta_t \bB_t x_t, \\  \nonumber
    y_t   &= \bC_t^{\top}\vh_t,
\end{align}
with state $\vh_t \in \mathbb{R}^N$, projections 
\[
\bB_t \coloneqq \begin{bmatrix} \bB_t \\ \hat{\bB}_t \end{bmatrix} \in \mathbb{R}^N, 
\qquad 
\bC_t \coloneqq \begin{bmatrix} \bC_t \\ -\hat{\bC}_t \end{bmatrix} \in \mathbb{R}^N, 
\]
and a transition matrix 
\[
\bR_t \;\coloneqq\; \text{Block}\Big(\{R(\Delta_t \boldsymbol{\theta_t}[i])\}_{i=1}^{N/2}\Big)
\in \mathbb{R}^{N \times N}
,
\qquad 
R(\theta) \coloneqq 
\begin{bmatrix}
    \cos(\theta) & -\sin(\theta) \\
    \sin(\theta) & \cos(\theta)
\end{bmatrix}.
\]
\end{restatable}

The proof is given in \cref{app:proof-complex-to-real-ssm}. 

\Cref{prop:complex-to-real-ssm} shows that the discretized complex SSM of state dimension $N/2$ has an equivalent real SSM with doubled state dimension ($N$), and its transition matrix is a scalar decayed block-diagonal matrix of $2\times2$ data-dependent rotation matrices ($e^{\Delta_t A_t} \bR_t$).

\begin{restatable}[Complex SSM, Data-Dependent RoPE Equivalence]{proposition}{PropRoPETrick}\label{prop:rope-trick} 
Under the notation established in \Cref{prop:complex-to-real-ssm}, consider the real SSM defined in \eqref{eq:real-ssm} unrolled for $T$ time-steps.
The output of the above SSM is equivalent to that of a vanilla scalar transition matrix-based SSM~\plaineqref{eq:mamba2-recurrence} with a data-dependent rotary embedding applied on the $\bB,\bC$ components of the SSM, as defined by:
\begin{equation}
    \vh_t = e^{\Delta_tA_t}\vh_{t-1} + \left( \prod_{i=0}^t \bR^\top_i \right) \Delta_t\bB_tx_t, \qquad \qquad
    y_t = \left[\left( \prod_{i=0}^t \bR^\top_i \right)\bC_t\right]^\top\vh_t \qquad
\end{equation}
where the matrix product represents right matrix multiplication, e.g., $\prod_{i=0}^1\bR_i = \bR_0 \bR_1$. 
We refer to the usage of a transformed real-valued SSM to compute the complex SSM as the ``RoPE trick.''
\end{restatable}
The proof is given in \cref{app:proof-rope-trick}.

To observe the connection of complex SSMs to RoPE embeddings, note that in the above proposition, the data-dependent rotations $\bR_i$ are aggregated across time-steps and applied to $\bC,\bB$, which, by the state space duality framework, correspond to the query ($\bQ$) and key ($\bK$) components of attention (\Cref{sec:prelim-sma}). 
Analogously, vanilla RoPE~\citep{su2023roformerenhancedtransformerrotary} applies \emph{data-independent} rotation matrices, where the rotation angles follow a fixed frequency schedule $\boldsymbol{\theta}[i] = 10000^{-2i/N}$. 

\subsubsection{Complex SSM with Exponential-Trapezoidal Discretization}

After deriving the recurrence for complex SSMs with exponential-Euler discretization, the generalization to exponential-trapezoidal discretization is similar.
\Cref{prop:rope-trick-trap} provides the full recurrence with the RoPE trick for Mamba-3.

\begin{restatable}[Rotary Embedding Equivalence with Exponential-Trapezoidal Discretization]{proposition}{PropRoPETrickTrap}\label{prop:rope-trick-trap}

Discretizing a complex SSM with the exponential-trapezoidal rule (\Cref{prop:trap}) yields the recurrence
\begin{align}
    \vh_t 
    &= \alpha_t \vh_{t-1} 
    + \beta_t \left(\prod_{i=0}^{t-1} \bR_i^\top \right)\bB_{t-1} x_{t-1} 
    + \gamma_t \left(\prod_{i=0}^t \bR_i^\top \right)\bB_t x_t, \nonumber \\
    y_t
    &= \left[\left(\prod_{i=0}^t \bR_i^\top\right)\bC_t\right]^\top \vh_t .
\end{align}
Here, $\bR_t$ is the block-diagonal rotation matrix defined in~\Cref{prop:complex-to-real-ssm}.
\end{restatable}

The proof is in \cref{app:proof-rope-trick-trap}.

We empirically validate that our complex SSM, implemented via data-dependent RoPE, is capable of solving state-tracking tasks that real-valued SSMs with and without standard RoPE cannot (\Cref{tab:reasoning}), supporting theoretical claims.
